\documentclass[11pt]{article}

\usepackage[margin=1in]{geometry}
\usepackage{array}
\usepackage{booktabs}
\usepackage{graphicx}
\usepackage{xurl}
\usepackage{hyperref}
\usepackage{xcolor}
\usepackage{enumitem}
\usepackage{float}
\usepackage{caption}
\usepackage{longtable}

\hypersetup{
  colorlinks=true,
  linkcolor=blue,
  urlcolor=blue,
  citecolor=blue
}

\setlist[itemize]{leftmargin=1.25em}
\setlist[enumerate]{leftmargin=1.5em}
\newcommand{\code}[1]{\texttt{\small #1}}
\newcommand{\filepath}[1]{\begingroup\urlstyle{tt}\small\url{#1}\endgroup}

\title{\textbf{Batch Effect Diagnostics for Taxonomy-aware BiomeGPT}}
\author{Yuanshu / Codex working report}
\date{May 20, 2026}

\begin{document}
\maketitle

\section*{Executive Summary}

We evaluated whether the taxonomy-aware BiomeGPT phase-2 sample embeddings encode
batch/study information and whether scGPT-style adversarial batch correction is a
reasonable next direction. The short answer is:

\begin{itemize}
  \item The batch labels used here are not raw metadata columns. They come from an
  externally enriched study/cohort batch-proxy annotation built from public accession
  lookup and conservative local-prefix rules.
  \item After receiving the uploaded real study-id table
  \code{BiomeGPT\_species\_samples\_studyIDs.csv}, we reimplemented the batch
  diagnostics and correction trials using true sample-to-study alignment rather than
  proxy labels.
  \item The phase-2 sample embeddings do contain measurable batch signal.
  \item However, most broad batch labels are strongly phenotype-confounded, so using all
  labels for adversarial correction would risk removing true disease biology.
  \item A conservative proof-of-concept using only safer batch labels ran end-to-end and
  reduced safe-batch predictability modestly while mostly preserving Healthy-vs-Diseased
  signal.
  \item This is a working diagnostic and correction pipeline, but not yet a final
  batch-corrected model.
\end{itemize}

The most important numerical result is shown below:

\begin{center}
\begin{tabular}{lcccc}
\toprule
Probe & Before macro-F1 & After macro-F1 & Before bal. acc. & After bal. acc. \\
\midrule
Safe batch labels & 0.499 & 0.476 & 0.503 & 0.475 \\
Healthy vs Diseased & 0.796 & 0.789 & 0.801 & 0.794 \\
\bottomrule
\end{tabular}
\end{center}

This means the light adversarial correction moved in the desired direction: batch
predictability decreased, while the biological Healthy-vs-Diseased signal was largely
preserved.

The newer real-study-id experiments are more stringent and should now be treated as the
main evidence. With true study ids, study signal is strong, phenotype confounding is
substantial, and naive GRL adversarial fine-tuning did not yet remove study signal. A
post-hoc representation correction can remove study signal almost completely, but broad
study correction also removes much of the H/D signal, confirming that study and phenotype
are entangled in this dataset.

\section{Scientific Question}

After building the taxonomy-aware BiomeGPT model, the next question was whether
the learned sample-level representations contain unwanted technical or cohort-specific
batch information. This matters because external validation performance may be limited
by study/cohort shift, and scGPT-style batch integration suggests a possible solution:
make sample embeddings less predictive of batch while retaining biological phenotype
information.

The goal of this stage was therefore not to immediately claim batch correction, but to
answer three diagnostic questions:

\begin{enumerate}
  \item Can batch/study labels be predicted from the BiomeGPT sample prompt?
  \item Are those batch labels safe to remove, or are they confounded with phenotype?
  \item If we run a conservative adversarial correction, does batch predictability go down
  without destroying Healthy-vs-Diseased signal?
\end{enumerate}

\section{Batch Annotation Construction}

\subsection{Why the data had to be processed}

The original phase-2 gut metadata did not contain an explicit technical batch column.
It also did not include a clean study ID, cohort ID, sequencing center, sequencing
platform, country, library preparation, or run-date column. The available fields were:

\begin{quote}
\code{sample\_id}, \code{Phenotype}, \code{body\_site}, \code{Phenotype\_fullname}
\end{quote}

Because adversarial batch correction requires a batch or domain label, we first had to
reconstruct the best available batch proxy. The goal was not to claim that every
technical sequencing batch is known. Instead, the goal was to build an auditable
study/cohort-level proxy that is good enough for diagnostic probing and cautious enough
for conservative correction experiments.

\subsection{How sample IDs were interpreted}

The key observation was that many \code{sample\_id} values contain public database
accessions. These accessions can be queried against public metadata services to recover
study or cohort information. Each sample ID was first classified by pattern:

\begin{center}
\begin{tabular}{ll}
\toprule
Pattern & Interpreted type \\
\midrule
\code{SAMN...} & NCBI BioSample accession \\
\code{SAMEA...} & ENA/EBI BioSample accession \\
\code{SAMD...} & DDBJ/INSDC BioSample accession \\
\code{SRR...} & NCBI SRA run accession \\
\code{EGAR...} & EGA run accession \\
Other & Local or unknown sample ID \\
\bottomrule
\end{tabular}
\end{center}

For compound IDs such as \code{EGAR00001420100\_9002000001328080LL}, only the public
run accession part was used for external lookup, while the original full sample ID was
preserved.

\subsection{External lookup strategy}

Public accession-like IDs were mapped to study or cohort metadata using official public
services:

\begin{center}
\begin{tabular}{lp{0.56\textwidth}}
\toprule
Source & What it provided \\
\midrule
ENA Portal API & Sample-to-study metadata for \code{SAMN}, \code{SAMEA}, and \code{SAMD} IDs \\
NCBI SRA RunInfo & Run-to-BioProject/SRAStudy metadata for \code{SRR} IDs \\
NCBI BioSample E-utilities & Fallback for older \code{SAMN} records missing usable ENA study accessions \\
EGA Metadata API & Run-to-study metadata for \code{EGAR} IDs \\
\bottomrule
\end{tabular}
\end{center}

The recommended batch label was then constructed from the strongest available study or
cohort identifier. Examples include:

\begin{quote}
\code{study:PRJNA763023}, \code{study:PRJEB103248+PRJEB22893},
\code{dbgap:phs000228}, \code{ega\_study:EGAS00001001704}
\end{quote}

If a sample mapped to multiple study accessions, the accessions were joined with
\code{+}. This can happen when a public database links both an original project and a
derived metagenomic assembly project.

\subsection{Local IDs and confidence levels}

Local sample IDs without public accessions could not be externally verified from this
metadata alone. They were therefore handled cautiously. A prefix-derived label was
assigned only when the local naming pattern was structured and biologically
interpretable. Examples include \code{PRIMM...}, \code{M...}, \code{SID...},
\code{wHAXPI...}, and \code{FMT...}. These labels are useful for diagnostics and manual
review, but they are not treated as fully verified external batch labels.

Every sample received an external confidence label:

\begin{center}
\begin{tabular}{lrl}
\toprule
Confidence & Samples & Meaning \\
\midrule
High & 9,389 & Exact public accession lookup returned study/cohort metadata \\
Medium & 3,418 & Structured local prefix-derived cohort label \\
Low & 525 & Unresolved or ambiguous local sample ID pattern \\
\bottomrule
\end{tabular}
\end{center}

This distinction matters because the batch probe can use broad labels for exploratory
diagnostics, but final correction should not treat prefix-derived or unresolved labels as
equally reliable.

\subsection{Phenotype-confounding check}

A real study label is not automatically safe for batch correction. Many microbiome
studies are highly phenotype-specific. For example, one public study may contain only
Healthy controls, while another may contain only Parkinson's disease samples. If we train
an adversarial model to remove those study labels, the model may also remove true disease
biology.

Therefore, each proposed batch label was checked for phenotype composition:

\begin{quote}
\code{n\_samples}, \code{n\_phenotypes}, \code{top\_phenotype},
\code{top\_phenotype\_count}, \code{top\_phenotype\_fraction}
\end{quote}

The warning rule was:

\begin{quote}
\code{phenotype\_confounding\_warning=True} if
\code{top\_phenotype\_fraction >= 0.80} or \code{n\_phenotypes <= 1}
\end{quote}

This is intentionally strict. If a batch group is nearly one phenotype, correcting away
that batch may also correct away the phenotype signal we want the model to preserve.

\subsection{Conservative-safe label rule}

For final batch-correction experiments, a sample was marked:

\begin{quote}
\code{safe\_for\_final\_batch\_correction\_conservative=True}
\end{quote}

only if both conditions were satisfied:

\begin{quote}
\code{external\_confidence == "high"} and
\code{phenotype\_confounding\_warning == False}
\end{quote}

This produced:

\begin{center}
\begin{tabular}{lr}
\toprule
Category & Samples \\
\midrule
Conservative-safe for final batch correction & 2,134 \\
Not conservative-safe & 11,198 \\
\bottomrule
\end{tabular}
\end{center}

The resulting annotation should therefore be described as an externally enriched
study/cohort batch-proxy annotation, not as a perfect technical sequencing-batch table.
The broad \code{batch\_label\_external\_recommended} field is appropriate for asking
whether embeddings contain study/cohort signal. The conservative-safe subset is the
appropriate starting point for asking whether adversarial correction can reduce batch
signal without destroying phenotype signal.

\section{Inputs Used}

\subsection{Model checkpoint}

We used the taxonomy-aware phase-2 checkpoint:

\begin{quote}
\code{dataset\_v3/outputs\_taxonomy\_notebook/taxonomy\_checkpoint\_stage2.pt}
\end{quote}

This checkpoint is the full taxonomy-aware run, not the earlier minimal BioGPT
checkpoint. Its model structure contains rank-wise taxonomy embeddings, including
Domain, Kingdom, Phylum, Class, Order, Family, Genus, and Species embeddings.

\subsection{Data}

The diagnostic data came from the phase-2 gut pretraining set and the externally enriched
batch annotation file:

\begin{quote}
\code{dataset\_v3/abund\_pretraining\_phase2\_gut.csv.zip}

\code{dataset\_v3/meta\_pretraining\_phase2\_gut\_batch\_annotation\_external\_enriched.csv}
\end{quote}

The final diagnostic matrix contained:

\begin{center}
\begin{tabular}{lr}
\toprule
Quantity & Count \\
\midrule
Phase-2 gut samples & 13,332 \\
Sample embedding dimension & 512 \\
Checkpoint species vocabulary & 1,012 \\
Unique recommended batch labels & 218 \\
High-confidence external labels & 9,389 \\
Medium-confidence labels & 3,418 \\
Low-confidence labels & 525 \\
Conservative-safe batch-correction samples & 2,134 \\
\bottomrule
\end{tabular}
\end{center}

\section{Method: Sample Prompt Extraction}

For each phase-2 gut sample, we extracted the BiomeGPT sample prompt, defined as the
final-layer \code{<cls>} embedding from the pretrained taxonomy-aware transformer.
The extraction procedure was:

\begin{enumerate}
  \item Load the taxonomy-aware phase-2 checkpoint.
  \item Load the phase-2 abundance matrix.
  \item Align abundance rows to the batch annotation metadata by sample ID.
  \item Reindex abundance columns to the checkpoint species vocabulary.
  \item Convert relative abundance values to 32 rank-based abundance bins.
  \item Run the pretrained model in inference mode and save the \code{<cls>} embedding.
\end{enumerate}

This produced an embedding matrix:

\begin{quote}
\code{X\_sample \(\in\) R\^{}(13332 x 512)}
\end{quote}

The cached output is:

\begin{quote}
\code{dataset\_v3/outputs\_batch\_diagnostics\_taxonomy/phase2\_sample\_prompt\_embeddings.npz}
\end{quote}

\section{Method: Batch Probe}

\subsection{What is a batch probe?}

A batch probe is a shallow supervised classifier trained on frozen model embeddings.
Here, the input is the sample prompt \code{<cls>} embedding, and the target is a
batch/study label. If the classifier predicts batch labels much better than a dummy
baseline, then the embedding contains batch information.

Importantly, this does not modify the BiomeGPT model. It is a diagnostic measurement.

\subsection{Probe model}

For each probe, we used:

\begin{itemize}
  \item Standard scaling of the 512-dimensional embeddings.
  \item Multiclass logistic regression.
  \item Class-balanced weighting.
  \item A stratified 80/20 train-test split when feasible.
  \item Metrics: accuracy, balanced accuracy, and macro-F1.
  \item Dummy baseline: most-frequent-class classifier.
\end{itemize}

Balanced accuracy and macro-F1 are more informative than raw accuracy because batch and
phenotype labels are imbalanced.

\subsection{Batch label panels}

We evaluated three batch-label panels:

\begin{enumerate}
  \item \textbf{High + medium labels}: broadest usable set for hypothesis generation.
  \item \textbf{High-only labels}: exact external accession lookups where available.
  \item \textbf{Conservative-safe labels}: labels passing the conservative safety rule for
  final batch correction.
\end{enumerate}

For each panel, labels with fewer than 10 samples were excluded from the probe to avoid
unstable tiny-class estimates.

\section{Batch Probe Results}

\begin{center}
\begin{tabular}{lrrrrr}
\toprule
Probe target & Samples & Classes & Bal. acc. & Macro-F1 & Dummy bal. acc. \\
\midrule
High + medium batch labels & 12,758 & 105 & 0.397 & 0.365 & 0.010 \\
High-only batch labels & 9,340 & 83 & 0.386 & 0.355 & 0.012 \\
Conservative-safe batch labels & 2,134 & 17 & 0.503 & 0.499 & 0.059 \\
Phenotype labels & 13,332 & 30 & 0.547 & 0.395 & 0.033 \\
Healthy vs Diseased & 13,332 & 2 & 0.801 & 0.796 & 0.500 \\
\bottomrule
\end{tabular}
\end{center}

\subsection{Interpretation}

The batch signal is real. For conservative-safe labels, the batch probe reached macro-F1
0.499, compared with dummy macro-F1 0.014. This is much higher than chance and shows
that study/batch information is encoded in the taxonomy-aware phase-2 sample prompts.

At the same time, the embeddings also retain strong biological signal: a shallow
Healthy-vs-Diseased probe reached macro-F1 0.796. This is important because batch
correction should reduce batch signal without removing phenotype signal.

\section{Method: Batch-Phenotype Confounding Analysis}

Before doing any correction, we checked whether batch labels were confounded with
phenotype. This is critical because if a batch label is nearly identical to a disease label,
then adversarially removing batch information may also remove real disease biology.

For each batch-label panel, we computed:

\begin{itemize}
  \item Batch-by-phenotype contingency table.
  \item Cramer's V between batch and phenotype.
  \item Normalized mutual information (NMI).
  \item Top phenotype fraction per batch.
  \item Number of batches where the top phenotype fraction is at least 0.8.
\end{itemize}

The top phenotype fraction is especially interpretable. For example, if one batch contains
100 samples and 95 are Healthy, its top phenotype fraction is 0.95. Correcting away that
batch could also correct away Healthy-specific biology.

\section{Confounding Results}

\begin{center}
\begin{tabular}{lrrrr}
\toprule
Label set & Cramer's V & NMI & Median top phenotype frac. & Batches top frac. $\geq$ 0.8 \\
\midrule
High + medium & 0.802 & 0.519 & 1.000 & 79 / 105 \\
High-only & 0.907 & 0.489 & 1.000 & 66 / 83 \\
Conservative-safe & 0.683 & 0.457 & 0.549 & 0 / 17 \\
\bottomrule
\end{tabular}
\end{center}

\subsection{Interpretation}

The broad high+medium and high-only batch labels are strongly phenotype-confounded.
Their median top phenotype fraction is 1.000, meaning many batches contain only one
phenotype. These labels are not safe for direct adversarial removal.

The conservative-safe label set is not perfectly independent of phenotype, but it is much
less confounded. Its median top phenotype fraction is 0.549, and none of its 17 retained
batches has top phenotype fraction above 0.8. Therefore, this set is the best starting
point for a cautious batch-correction proof-of-concept.

\section{Conservative Adversarial Correction Trial}

\subsection{Why adversarial correction?}

scGPT-style batch integration uses a domain-adversarial branch: a discriminator tries to
predict batch from the sample representation, while gradient reversal makes the encoder
learn representations that are less batch-predictive. We adapted that idea to BiomeGPT by
using the \code{<cls>} sample prompt as the representation for batch discrimination.

\subsection{What exactly was corrected?}

This was not a new pretraining run from scratch. The trial continued training from the
existing taxonomy-aware phase-2 checkpoint:

\begin{quote}
\filepath{dataset_v3/outputs_taxonomy_notebook/taxonomy_checkpoint_stage2.pt}
\end{quote}

The correction target was the sample-level \code{<cls>} embedding produced by the
taxonomy-aware transformer. This is the same sample prompt used by the downstream probes.
The goal was to make this \code{<cls>} embedding less predictive of conservative-safe
study/cohort labels while keeping the original masked-abundance reconstruction behavior.

We did not train on all batch labels. The adversarial branch used only rows satisfying:

\begin{quote}
\code{safe\_for\_final\_batch\_correction\_conservative=True}
\end{quote}

and retained only batch classes with at least 10 samples. This gave:

\begin{center}
\begin{tabular}{lr}
\toprule
Quantity & Value \\
\midrule
Safe training samples & 2,134 \\
Batch classes & 17 \\
Minimum retained class size & 43 \\
Maximum retained class size & 278 \\
\bottomrule
\end{tabular}
\end{center}

\subsection{Implementation details}

The implementation is in:

\begin{quote}
\filepath{dataset_v3/batch_adversarial_correction.py}
\end{quote}

The data flow was:

\begin{enumerate}
  \item Load the taxonomy-aware phase-2 checkpoint.
  \item Load the phase-2 gut abundance matrix.
  \item Align abundance rows to the enriched batch annotation table.
  \item Filter to conservative-safe samples.
  \item Reindex abundance columns to the checkpoint species vocabulary.
  \item Convert abundance values to 32 rank-based bins.
  \item Continue training the BiomeGPT encoder with the original reconstruction objective
  plus a small adversarial batch objective on the \code{<cls>} token.
\end{enumerate}

\subsection{Batch discriminator}

The adversarial branch was a small multilayer perceptron attached to the \code{<cls>}
embedding. Its architecture was:

\begin{quote}
\code{LayerNorm(512) -> Linear(512, 256) -> ReLU -> Dropout(0.1)}

\code{-> Linear(256, 128) -> ReLU -> Dropout(0.1) -> Linear(128, 17)}
\end{quote}

The discriminator's direct task is ordinary multiclass batch classification. It receives
the \code{<cls>} embedding and predicts one of the 17 conservative-safe batch labels.

\subsection{Gradient reversal mechanism}

The important point is that the discriminator is connected through a gradient reversal
layer (GRL). In the forward pass, GRL behaves like the identity function:

\begin{quote}
\code{GRL(x) = x}
\end{quote}

In the backward pass, it multiplies the gradient by \code{-lambda}. In this trial:

\begin{quote}
\code{grl\_lambda = 1.0}
\end{quote}

This creates an adversarial objective:

\begin{itemize}
  \item The discriminator learns to predict batch from the \code{<cls>} embedding.
  \item The encoder receives the reversed gradient, so it learns to make the
  \code{<cls>} embedding less useful for batch prediction.
\end{itemize}

Therefore, the discriminator and encoder have opposite goals. The discriminator tries to
recover batch labels; the encoder tries to hide batch information while still solving the
masked abundance reconstruction task.

\subsection{Training objective}

For each mini-batch, we generated a random species-token mask using:

\begin{quote}
\code{mask\_ratio = 0.25}
\end{quote}

The model then encoded the binned abundance vector with the mask applied. The species
token hidden states were sent to the reconstruction head, and the \code{<cls>} hidden
state was sent to the batch discriminator.

The optimized loss was:

\begin{quote}
\code{total loss = reconstruction loss + dab\_weight * batch cross-entropy loss}
\end{quote}

with:

\begin{quote}
\code{dab\_weight = 0.1}
\end{quote}

The reconstruction loss was mean-squared error on masked nonzero abundance bins. The
batch loss was cross-entropy over 17 conservative-safe batch classes. Because the batch
loss passes through GRL, minimizing this total loss trains the discriminator normally but
pushes the encoder in the opposite direction for the batch component.

\subsection{Why this is a light correction}

This run was intentionally conservative. The settings were:

\begin{center}
\begin{tabular}{lr}
\toprule
Hyperparameter & Value \\
\midrule
Epochs & 3 \\
Training batch size & 8 \\
Learning rate & 2e-5 \\
Weight decay & 1e-4 \\
Mask ratio & 0.25 \\
DAB weight & 0.1 \\
GRL lambda & 1.0 \\
Discriminator dropout & 0.1 \\
\bottomrule
\end{tabular}
\end{center}

It is ``light'' because the model is only continued for 3 epochs and the adversarial term
has small weight. The main training pressure remains the original masked-abundance
reconstruction objective. This reduces the risk of damaging the pretrained microbiome
representation before we have downstream validation.

\subsection{Checkpoint and evaluation workflow}

After the 3-epoch correction run, the corrected model was saved as:

\begin{quote}
\filepath{dataset_v3/outputs_batch_adversarial_taxonomy/taxonomy_checkpoint_stage2_batch_adversarial_safe.pt}
\end{quote}

Then we re-extracted \code{<cls>} embeddings for all 13,332 phase-2 gut samples from the
corrected checkpoint and reran the same frozen logistic-regression probes:

\begin{itemize}
  \item safe-batch probe on conservative-safe batch labels;
  \item Healthy-vs-Diseased probe on all samples.
\end{itemize}

The desired outcome was not maximum downstream disease performance yet. The immediate
criterion was directional: safe-batch probe performance should decrease, while
Healthy-vs-Diseased probe performance should remain close to the uncorrected baseline.

\section{Adversarial Trial Results}

\subsection{Trial A: light correction}

Settings:

\begin{itemize}
  \item Epochs: 3
  \item DAB weight: 0.1
  \item Gradient reversal lambda: 1.0
  \item Safe training samples: 2,134
  \item Batch classes: 17
\end{itemize}

Training history:

\begin{center}
\begin{tabular}{lrrrr}
\toprule
Epoch & Total loss & Reconstruction loss & DAB loss & DAB train acc. \\
\midrule
1 & 37.924 & 37.645 & 2.790 & 0.103 \\
2 & 37.471 & 37.198 & 2.733 & 0.143 \\
3 & 37.607 & 37.337 & 2.703 & 0.147 \\
\bottomrule
\end{tabular}
\end{center}

The DAB training accuracy increased from 0.103 to 0.147, showing that the discriminator
was learning some batch signal during training. Because the discriminator was connected
through GRL, that same signal produced reversed gradients for the encoder.

Before/after probe results:

\begin{center}
\begin{tabular}{lrrrr}
\toprule
Probe & Before macro-F1 & After macro-F1 & Before bal. acc. & After bal. acc. \\
\midrule
Safe batch labels & 0.499 & 0.476 & 0.503 & 0.475 \\
Healthy vs Diseased & 0.796 & 0.789 & 0.801 & 0.794 \\
\bottomrule
\end{tabular}
\end{center}

This is the best current trial. Batch predictability decreased modestly, while
Healthy-vs-Diseased predictability was mostly preserved.

\subsection{Trial B: stronger correction}

Settings:

\begin{itemize}
  \item Epochs: 5
  \item DAB weight: 0.5
  \item Gradient reversal lambda: 1.0
\end{itemize}

\begin{center}
\begin{tabular}{lrrrr}
\toprule
Probe & Before macro-F1 & After macro-F1 & Before bal. acc. & After bal. acc. \\
\midrule
Safe batch labels & 0.499 & 0.505 & 0.503 & 0.505 \\
Healthy vs Diseased & 0.796 & 0.786 & 0.801 & 0.790 \\
\bottomrule
\end{tabular}
\end{center}

This stronger run did not improve batch correction. Batch predictability slightly
increased, and H/D signal decreased slightly. This suggests that simply increasing
adversarial weight is not sufficient; the schedule and training design need tuning.

\section{Additional Label-Panel Ablation}

After the initial conservative-safe run, we tested three label-panel choices under the
same light adversarial setting:

\begin{enumerate}
  \item \textbf{High + medium labels}: all exact external labels plus structured
  prefix-derived labels.
  \item \textbf{High-only labels}: exact external database labels only, without removing
  phenotype-confounded batches.
  \item \textbf{Conservative-safe labels}: exact external labels after removing strongly
  phenotype-confounded batches.
\end{enumerate}

The purpose was to separate two questions:

\begin{itemize}
  \item Does using broader labels create stronger batch removal?
  \item Does increasing training scale or strength necessarily improve correction?
\end{itemize}

For this ablation, the script was generalized with:

\begin{quote}
\code{--label\_panel high\_medium | high\_only | conservative\_safe}
\end{quote}

and the three runs used the same light correction hyperparameters:

\begin{quote}
\code{epochs=3, dab\_weight=0.1, grl\_lambda=1.0, train\_batch\_size=32}
\end{quote}

\subsection{Ablation results}

\begin{center}
\begin{tabular}{lrrrr}
\toprule
Training panel & Samples & Classes & Batch macro-F1 before & Batch macro-F1 after \\
\midrule
High + medium & 12,758 & 105 & 0.365 & 0.383 \\
High-only & 9,340 & 83 & 0.355 & 0.373 \\
Conservative-safe & 2,134 & 17 & 0.499 & 0.511 \\
\bottomrule
\end{tabular}
\end{center}

\begin{center}
\begin{tabular}{lrrr}
\toprule
Training panel & Batch bal. acc. before & Batch bal. acc. after & H/D macro-F1 after \\
\midrule
High + medium & 0.397 & 0.411 & 0.804 \\
High-only & 0.386 & 0.405 & 0.802 \\
Conservative-safe & 0.503 & 0.513 & 0.794 \\
\bottomrule
\end{tabular}
\end{center}

In this batch-size-32 ablation, none of the three label-panel runs reduced the target
batch probe. The high+medium and high-only runs actually increased batch predictability.
This is not the desired correction direction.

\subsection{Comparison against earlier runs}

The best directional result is still the earlier conservative-safe light run with
\code{train\_batch\_size=8}:

\begin{center}
\begin{tabular}{lrrrr}
\toprule
Run & Batch macro-F1 before & Batch macro-F1 after & H/D macro-F1 before & H/D macro-F1 after \\
\midrule
Conservative-safe, batch size 8, DAB 0.1, 3 epochs & 0.499 & 0.476 & 0.796 & 0.789 \\
Conservative-safe, batch size 32, DAB 0.1, 3 epochs & 0.499 & 0.511 & 0.796 & 0.794 \\
Conservative-safe, DAB 0.5, 5 epochs & 0.499 & 0.505 & 0.796 & 0.786 \\
\bottomrule
\end{tabular}
\end{center}

This answers the ``should we simply train more or larger?'' question: probably not by
itself. Stronger or larger updates did not reliably improve batch removal. In fact, the
broader label panels and the stronger DAB run moved in the wrong direction for the batch
probe.

\subsection{Interpretation}

These ablations support three conclusions:

\begin{itemize}
  \item Broad high+medium labels are useful for diagnostics, but they are too heterogeneous
  and phenotype-confounded for naive adversarial correction.
  \item High-only labels are externally verified, but many are still phenotype-confounded;
  external verification alone does not make a label safe for correction.
  \item Conservative-safe labels remain the most scientifically defensible correction
  target, but optimization is sensitive. We should tune schedule and validation criteria
  rather than only increasing epochs or DAB weight.
\end{itemize}

The next technical step should be a small hyperparameter sweep on conservative-safe labels:
for example, DAB weight \code{0.02, 0.05, 0.1}, batch size \code{8 or 16}, and a gradual
GRL/DAB warm-up schedule. The selection criterion should require lower conservative-safe
batch probe performance while preserving Healthy-vs-Diseased probe performance.

\section{Masked Abundance Reconstruction Diagnostic}

\subsection{Why this diagnostic matters}

BiomeGPT is a foundation model whose pretraining objective is masked abundance prediction.
Healthy-vs-Diseased classification is a downstream task. Therefore, batch correction should
not be evaluated only by H/D probe performance. A better foundation-model criterion is:

\begin{quote}
Does correction reduce batch predictability while preserving or improving masked
abundance reconstruction?
\end{quote}

To evaluate this directly, we added:

\begin{quote}
\filepath{dataset_v3/batch_reconstruction_diagnostics.py}
\end{quote}

For each checkpoint, the script applies the same deterministic mask to nonzero abundance
bins and computes reconstruction MSE on the masked positions. It also computes
batch-stratified MSE summaries for the selected label panels.

\subsection{Initial reconstruction results}

\begin{center}
\small
\begin{tabular}{p{0.34\textwidth}rrrr}
\toprule
Checkpoint & Overall MSE & Safe-panel MSE & Safe-panel CV & Safe-panel gap \\
\midrule
Original phase-2 & 34.857 & 37.808 & 0.095 & 13.140 \\
Conservative-safe, batch size 8 & 33.745 & 35.724 & 0.093 & 12.961 \\
Conservative-safe, batch size 32 & 34.088 & 36.252 & 0.096 & 13.324 \\
High + medium, batch size 32 & 33.049 & 35.770 & 0.096 & 12.736 \\
High-only, batch size 32 & 33.373 & 36.061 & 0.098 & 13.453 \\
Conservative-safe, DAB 0.5 & 33.853 & 35.256 & 0.092 & 12.554 \\
\bottomrule
\end{tabular}
\end{center}

The main observation is that continued correction training did not destroy the masked
abundance objective. Overall reconstruction MSE decreased for all corrected checkpoints
relative to the original phase-2 checkpoint. The best conservative-safe batch-size-8 run
also reduced safe-panel reconstruction MSE and slightly reduced the safe-panel
batch-MSE gap.

\subsection{Interpretation}

This improves the evaluation design. The correct selection criterion should have three
layers:

\begin{enumerate}
  \item \textbf{Batch removal}: post-hoc batch probe should decrease.
  \item \textbf{Foundation objective preservation}: masked abundance reconstruction MSE
  should remain stable or improve, ideally with lower per-batch MSE gap/CV.
  \item \textbf{Downstream biological sanity}: H/D probe or ExVal performance should not
  collapse.
\end{enumerate}

Importantly, lower reconstruction MSE alone is not enough. Some runs improved
reconstruction while the batch probe increased. This means abundance modeling and batch
removal must be evaluated jointly.

\section{Discussion and Scientific Insights}

\subsection{Batch labels are real, but not all are safe correction targets}

The batch probe results show that the taxonomy-aware BiomeGPT sample embeddings encode
study/cohort-associated signal. This supports the motivation for batch-effect diagnostics.
However, the confounding analysis shows that many candidate batch labels are still tied to
phenotype. This is especially true for the broad high+medium and high-only panels.

The main insight is:

\begin{quote}
External study labels are useful for diagnostics, but external verification alone does not
make them safe for correction.
\end{quote}

A label can be a real public study label and still be unsafe for adversarial correction if
the study contains mostly one disease phenotype.

\subsection{Phenotype confounding likely explains broad-label failure}

For both broad-label adversarial trials, the target batch probe increased after training:

\begin{center}
\begin{tabular}{lrr}
\toprule
Panel & Batch macro-F1 before & Batch macro-F1 after \\
\midrule
High + medium & 0.365 & 0.383 \\
High-only & 0.355 & 0.373 \\
\bottomrule
\end{tabular}
\end{center}

This suggests that the model did not learn a clean technical-batch removal operation.
Because study and phenotype are entangled, the adversarial objective may push the encoder
into a new representation that still separates study/cohort structure, or even separates
phenotype-associated study structure more clearly.

\subsection{More aggressive training is not automatically better}

The stronger conservative-safe run also did not improve batch removal:

\begin{center}
\begin{tabular}{lrr}
\toprule
Run & Batch macro-F1 after & H/D macro-F1 after \\
\midrule
Conservative-safe, DAB 0.1, batch size 8, 3 epochs & 0.476 & 0.789 \\
Conservative-safe, DAB 0.5, 5 epochs & 0.505 & 0.786 \\
\bottomrule
\end{tabular}
\end{center}

Therefore, the current limitation is probably not simply that the model was trained too
little. Increasing DAB weight or epochs can make the adversarial objective less stable and
may damage biological signal without removing batch signal.

\subsection{The discriminator-encoder game may be unstable}

Adversarial correction is a two-player optimization problem. The discriminator tries to
predict batch from \code{<cls>}; the encoder tries to hide batch from the discriminator;
the reconstruction loss tries to preserve microbiome abundance structure. If the
discriminator learns too quickly, too slowly, or from noisy/confounded labels, the reversed
gradient may not point toward a clean batch-invariant representation.

This instability is a plausible explanation for why several runs increased post-hoc batch
probe performance instead of decreasing it.

\subsection{The encoder may learn new batch-separable structure}

The adversarial discriminator only penalizes the batch directions it can exploit during
training. After correction, we evaluate with a fresh frozen logistic-regression probe. If
the after-probe performs better, that means batch information is still present, possibly in
a different geometry. In other words, the encoder may hide batch from the training
discriminator while still producing embeddings that remain separable by another classifier.

This is why post-hoc probes are important: they test whether the final embedding space is
actually less batch-predictive, not merely whether the training discriminator struggled.

\subsection{No warm-up schedule is a weakness}

The current implementation applies fixed adversarial pressure from the beginning:

\begin{quote}
\code{grl\_lambda = 1.0}, \quad \code{dab\_weight = fixed}
\end{quote}

This may be too abrupt. A more stable domain-adversarial setup would warm up the
adversarial signal gradually:

\begin{quote}
start with low or zero DAB pressure, let reconstruction stabilize, then gradually increase
GRL/DAB strength.
\end{quote}

This could prevent early adversarial gradients from pushing the pretrained representation
into an unstable or newly batch-separable geometry.

\subsection{Batch size changes adversarial dynamics}

Batch size is not only a speed parameter in adversarial training. The best directional
result used conservative-safe labels with batch size 8:

\begin{quote}
\code{batch macro-F1: 0.499 -> 0.476}
\end{quote}

but the otherwise similar batch-size-32 conservative-safe run moved in the wrong
direction:

\begin{quote}
\code{batch macro-F1: 0.499 -> 0.511}
\end{quote}

This suggests that minibatch composition affects the discriminator-encoder game. Smaller
batches may regularize the adversarial dynamics differently, while larger batches may make
the discriminator updates smoother but not necessarily more useful for batch-invariant
representations.

\subsection{Overall interpretation}

The data support the batch-correction idea, but not a final corrected model yet. A concise
interpretation is:

\begin{quote}
Taxonomy-aware BiomeGPT embeddings encode study/cohort-associated signal. However, broad
batch labels are strongly phenotype-confounded, and naive adversarial correction can fail
or even increase post-hoc batch separability. The most promising direction is
conservative-safe labels with a more careful adversarial schedule, rather than simply using
more labels, more epochs, or a stronger DAB loss.
\end{quote}

\section{Real Study-ID Reimplementation}

\subsection{Why this was necessary}

The first batch pipeline used externally enriched proxy batch labels because the raw
phase-2 metadata did not directly expose a clean study-id field. After the file
\filepath{BiomeGPT_species_samples_studyIDs.csv} became available, we rebuilt the batch
annotation around true sample-to-study mapping. This is scientifically better because the
batch variable is now a real cohort/study identifier rather than an inferred proxy.

The uploaded CSV had sample ids in the first column and study names in \code{study\_name}.
Some rows contained unmatched quote characters, so the parser reads the file with quoting
disabled, then strips quotes and normalizes whitespace. We then align by exact
\code{sample\_id} to \filepath{dataset_v3/meta_pretraining_phase2_gut.csv}.

\subsection{Real study-id coverage}

\begin{center}
\begin{tabular}{lr}
\toprule
Quantity & Count \\
\midrule
Rows in uploaded study-id file & 13,524 \\
Phase-2 gut samples & 13,332 \\
Phase-2 samples matched to real study id & 13,326 \\
Phase-2 samples missing real study id & 6 \\
Unique real studies in phase-2 gut & 80 \\
Conservative-safe samples & 4,751 \\
Conservative-safe studies & 26 \\
\bottomrule
\end{tabular}
\end{center}

For matched samples, the new batch label is:

\begin{quote}
\code{batch\_label\_external\_recommended = real\_study:<study\_name>}
\end{quote}

All matched real-study labels are treated as high-confidence because they come from the
uploaded explicit sample-to-study table. The old proxy labels remain useful history, but
the real-study results supersede them for the next round of experiments.

\subsection{Real study-id diagnostic probes}

Using the original taxonomy-aware phase-2 checkpoint, we re-extracted the 512-dimensional
sample-prompt embeddings and trained the same shallow logistic probes. Results:

\begin{center}
\begin{tabular}{lrrrr}
\toprule
Probe panel & Samples & Classes & Macro-F1 & Balanced acc. \\
\midrule
Real study, high-only & 13,326 & 80 & 0.456 & 0.495 \\
Real study, conservative-safe & 4,751 & 26 & 0.518 & 0.537 \\
Healthy vs Diseased & 13,332 & 2 & 0.796 & 0.801 \\
\bottomrule
\end{tabular}
\end{center}

The dummy balanced accuracy for 80 studies is only 0.0125, so a balanced accuracy near
0.495 means the pretrained sample embeddings contain substantial real study information.
The conservative-safe panel is not weaker; it is actually slightly easier to classify by
study, probably because it contains fewer but more coherent studies.

\subsection{Real study-id phenotype confounding}

The real study labels confirm the central risk: study id and phenotype are not
independent.

\begin{center}
\begin{tabular}{lrrrr}
\toprule
Panel & Studies & Cramer's V & NMI & Studies with top phenotype $\geq 0.8$ \\
\midrule
High-only real study & 80 & 0.792 & 0.526 & 54 \\
Conservative-safe real study & 26 & 0.634 & 0.528 & 0 \\
\bottomrule
\end{tabular}
\end{center}

For the full high-only real-study panel, 52 of 80 studies are phenotype-pure under the
available labels. This means that blindly removing all real-study information would also
remove disease-associated structure. The conservative-safe filter is therefore still
important: it does not remove all confounding, but it avoids the most obvious
single-phenotype studies.

\subsection{Real study-id adversarial correction trials}

We then reran GRL/DAB adversarial correction using true study ids. This implementation
continues phase-2 masked-abundance training and adds a domain-adversarial batch
classifier on the \code{<cls>} embedding. We also added warm-up parameters:

\begin{itemize}
  \item \code{dab\_warmup\_epochs}: gradually increases DAB loss weight.
  \item \code{grl\_warmup\_epochs}: gradually increases gradient-reversal strength.
\end{itemize}

The exact loss remains:

\begin{quote}
\code{total loss = masked-abundance reconstruction loss + effective\_dab\_weight * batch CE}
\end{quote}

with gradient reversal applied to the batch classifier input, so the discriminator learns
to predict study while the encoder receives the reversed gradient.

\begin{center}
\begin{tabular}{lrrrr}
\toprule
Trial & Study macro-F1 before & Study macro-F1 after & H/D before & H/D after \\
\midrule
Safe, DAB 0.05, batch 8, warm-up & 0.518 & 0.580 & 0.796 & 0.798 \\
Safe, DAB 0.10, batch 8, warm-up & 0.518 & 0.547 & 0.796 & 0.810 \\
High-only, DAB 0.05, batch 16, warm-up & 0.456 & 0.493 & 0.796 & 0.804 \\
\bottomrule
\end{tabular}
\end{center}

This is an important negative result. The model fine-tuning runs did not remove real
study signal; independent post-hoc probes found study signal increased. At the same time,
H/D performance stayed stable or improved. The most likely interpretation is that the
reconstruction objective plus weak adversarial pressure continued to organize the
embedding space around real cohort/phenotype structure, while the DAB objective was not
strong or stable enough to force study-invariant representations.

\subsection{Masked-abundance reconstruction after real-study GRL}

Because BiomeGPT is fundamentally a masked-abundance foundation model, we also evaluated
the main pretraining objective directly. We used the same deterministic nonzero-species
mask for all checkpoints and computed masked abundance-bin MSE.

\begin{center}
\begin{tabular}{lrrr}
\toprule
Checkpoint & Overall MSE & Safe-panel MSE & High-only MSE \\
\midrule
Base phase-2 & 34.856 & 36.221 & 34.857 \\
Safe GRL, DAB 0.05 & 33.235 & 33.852 & 33.235 \\
Safe GRL, DAB 0.10 & 33.164 & 33.755 & 33.165 \\
High-only GRL, DAB 0.05 & 32.911 & 34.225 & 32.911 \\
\bottomrule
\end{tabular}
\end{center}

So the GRL trials did preserve and even improve masked-abundance reconstruction. The
failure mode is not catastrophic forgetting. The failure mode is more specific: the
model remains, or becomes more, linearly study-separable in the \code{<cls>} embedding.

\subsection{Representation-level real-study correction baseline}

To test whether real study signal is removable at all, we implemented a transparent
post-hoc embedding correction baseline:

\begin{quote}
\code{x\_corrected = x - cross\_fitted\_study\_mean + cross\_fitted\_global\_mean}
\end{quote}

This is a mean-only study correction. We also tested a stronger mean+scale version that
maps each study's embedding variance to the global variance. The correction is
cross-fitted with stratified folds, so each sample is corrected using study means
estimated from other samples, not from itself.

\begin{center}
\begin{tabular}{llrrrr}
\toprule
Panel & Method & Study F1 before & Study F1 after & H/D F1 before & H/D F1 after \\
\midrule
High-only & Mean-center & 0.456 & 0.002 & 0.795 & 0.473 \\
High-only & Mean+scale & 0.456 & 0.008 & 0.795 & 0.471 \\
Conservative-safe & Mean-center & 0.518 & 0.004 & 0.690 & 0.637 \\
Conservative-safe & Mean+scale & 0.518 & 0.031 & 0.690 & 0.644 \\
\bottomrule
\end{tabular}
\end{center}

This result is very useful. It shows that study signal can be removed from embeddings
when we explicitly subtract study-level structure. However, it also shows why broad
study correction is dangerous: high-only correction nearly destroys H/D separability.
The conservative-safe panel has a much smaller biological cost, which supports the
earlier decision to avoid obviously phenotype-confounded study labels.

\subsection{Updated interpretation from real study ids}

The real-study-id experiments change the story from ``we have a promising light GRL
correction'' to a more precise conclusion:

\begin{quote}
Real study signal is strong and removable, but it is entangled with phenotype. Current
model-level GRL fine-tuning preserves reconstruction but does not yet remove real study
signal. Explicit embedding-level study centering removes study signal, and the biological
damage is acceptable only under conservative-safe filtering.
\end{quote}

The next model-level step should therefore not simply increase epochs. It should translate
the successful representation-level behavior into the training objective, for example by
adding a mean-alignment penalty such as study-centroid matching/CORAL/MMD on
\code{<cls>} embeddings while keeping masked-abundance reconstruction as the primary
loss, and selecting checkpoints by three criteria: lower real-study probe, preserved
masked-abundance MSE, and preserved H/D probe.

\subsection{Follow-up: representation alignment inside training}

We then tested whether the successful post-hoc real-study centering could be moved into
model training. Two model-level variants were implemented.

First, we added minibatch centroid alignment:

\begin{quote}
\code{loss = reconstruction loss + alignment\_weight * centroid\_alignment\_loss}
\end{quote}

The training loader was changed to use balanced study minibatches, so each minibatch
contains multiple studies and multiple samples per study. This makes centroid estimates
meaningful inside the batch.

Second, we implemented cross-fitted centroid distillation. We first compute the same
corrected targets that worked in post-hoc analysis:

\begin{quote}
\code{target = base\_embedding - cross\_fitted\_study\_mean + cross\_fitted\_global\_mean}
\end{quote}

Then the model is trained with:

\begin{quote}
\code{loss = reconstruction loss + target\_weight * MSE(unmasked\_<cls>, target)}
\end{quote}

An important debugging point: the target loss must be applied to the unmasked
\code{sample\_prompt} embedding, because that is the representation used by the probes.
Applying it to the masked pretraining embedding does not directly constrain the evaluated
embedding.

\begin{center}
\begin{tabular}{lrrrrr}
\toprule
Method & Study F1 before & Study F1 after & H/D before & H/D after & Recon MSE \\
\midrule
Base phase-2 & 0.518 & -- & 0.796 & -- & 34.856 \\
Centroid align w10 & 0.518 & 0.598 & 0.796 & 0.807 & 33.697 \\
Centroid align w50 & 0.518 & 0.579 & 0.796 & 0.819 & 33.628 \\
Unmasked distill w10 & 0.518 & 0.532 & 0.796 & 0.801 & 33.656 \\
Unmasked distill w50 & 0.518 & 0.556 & 0.796 & 0.802 & 33.587 \\
Unmasked distill w200 & 0.518 & 0.576 & 0.796 & 0.800 & 33.656 \\
\bottomrule
\end{tabular}
\end{center}

These are useful negative results. All model-level alignment trials preserve or improve
masked-abundance reconstruction, and H/D signal remains usable. However, none of these
unconstrained backbone fine-tuning runs reduce real-study predictability. The independent
study probe remains stronger than the original phase-2 embedding.

The practical result is therefore:

\begin{quote}
Direct cross-fitted study centering currently works as a corrected representation, but
simple GRL, minibatch centroid alignment, and centroid-target distillation are not yet
sufficient to make the BiomeGPT backbone itself study-invariant.
\end{quote}

The corrected embedding files were saved for downstream use:

\begin{itemize}
  \item \filepath{dataset_v3/outputs_real_study_embedding_correction_saved/real_study_high_only_mean_center_corrected_embeddings.npz}
  \item \filepath{dataset_v3/outputs_real_study_embedding_correction_saved/real_study_conservative_safe_mean_center_corrected_embeddings.npz}
\end{itemize}

\subsection{Follow-up: scGPT-style batch-conditioned decoder}

Because scGPT uses batch-aware mechanisms such as batch labels, batch embeddings, and
domain-specific batch normalization, we also tested whether BiomeGPT failed because it had
no batch side-channel. The hypothesis was:

\begin{quote}
If study id is available to the reconstruction decoder, reconstruction can model
study-specific abundance shifts without forcing the final sample embedding to carry all
study information.
\end{quote}

We implemented a batch-conditioned residual reconstruction decoder:

\begin{quote}
\code{pred = base\_reconstruction\_head(h\_species) + residual\_head(h\_species, study\_id)}
\end{quote}

The study id is not given to \code{sample\_prompt}; it is only used by the residual
decoder. The sample embedding is still trained toward the cross-fitted corrected target.
We also tested an optional DAB loss and a detached reconstruction branch where
reconstruction gradients do not update the encoder.

\begin{center}
\begin{tabular}{lrrrr}
\toprule
Trial & Study F1 before & Study F1 after & H/D before & H/D after \\
\midrule
Batch decoder, normalized target w50 & 0.518 & 0.565 & 0.796 & 0.805 \\
Batch decoder + DAB 0.1 & 0.518 & 0.577 & 0.796 & 0.804 \\
Batch decoder, detached recon encoder & 0.518 & 0.592 & 0.796 & 0.796 \\
Batch decoder, raw target w1 & 0.518 & 0.558 & 0.796 & 0.803 \\
\bottomrule
\end{tabular}
\end{center}

This did not solve real-study removal. The trials preserved H/D signal, but the
independent real-study probe remained stronger than the base phase-2 embedding. The
important interpretation is that adding a decoder-side batch residual only during late
continued training is probably not equivalent to scGPT-style batch-aware pretraining.
The batch-aware mechanism likely needs to be part of the architecture and training
dynamics from earlier in pretraining, or implemented as a frozen embedding correction
adapter rather than as unconstrained backbone fine-tuning.

\subsection{Follow-up: batch-token continued pretraining}

We then created a cleaner batch-aware continued-pretraining script:

\begin{quote}
\filepath{dataset_v3/batch_token_pretraining.py}
\end{quote}

During masked-abundance pretraining, the model sequence is:

\begin{quote}
\code{[CLS, STUDY, species\_1, species\_2, ...]}
\end{quote}

The attention mask lets species tokens attend to the \code{STUDY} token, so reconstruction
has a batch side-channel. However, \code{CLS} does not attend to \code{STUDY}, and
downstream embeddings are still extracted with the original unconditioned
\code{model.sample\_prompt(bins)}.

\begin{center}
\begin{tabular}{lrrrrr}
\toprule
Trial & Study F1 before & Study F1 after & H/D before & H/D after & Recon MSE \\
\midrule
Base phase-2 & 0.518 & -- & 0.796 & -- & 34.856 \\
Batch token, reconstruction only & 0.518 & 0.529 & 0.796 & 0.799 & 33.582 \\
Batch token + DAB 0.1 & 0.518 & 0.548 & 0.796 & 0.805 & 33.594 \\
Batch token + corrected-target distill & 0.518 & 0.544 & 0.796 & 0.800 & 33.584 \\
\bottomrule
\end{tabular}
\end{center}

This is the most stable model-level batch-aware design so far. It improves
masked-abundance reconstruction and keeps H/D signal, while the reconstruction-only run
increases study separability only mildly. However, it still does not reduce real-study
probe performance below the base phase-2 embedding. Adding DAB or corrected-target
distillation made the independent study probe stronger.

The interpretation is that a true study token side-channel helps reconstruction, but
adding it after the phase-2 representation is already formed is probably too late to make
the final sample embedding batch-invariant. To fully test the scGPT-style idea, the
batch-token mechanism should be present earlier in pretraining, ideally before or at the
start of phase2.

\subsection{Recommended experimental design}

The next experiment should be a controlled conservative-safe hyperparameter sweep:

\begin{itemize}
  \item Label panel: conservative-safe only.
  \item Batch size: \code{8} or \code{16}.
  \item DAB weight: \code{0.02}, \code{0.05}, \code{0.1}.
  \item GRL/DAB schedule: gradual warm-up over the first 30--50\% of epochs.
  \item Model selection: lower conservative-safe batch probe, preserved H/D probe, and no
  large reconstruction degradation.
\end{itemize}

\section{Files Produced}

\begin{longtable}{p{0.42\textwidth}p{0.50\textwidth}}
\toprule
File & Description \\
\midrule
\endhead
\filepath{dataset_v3/batch_effect_diagnostics.py} &
Script for extracting sample prompts, training batch/phenotype probes, computing
confounding, and saving PCA plots. \\
\filepath{dataset_v3/batch_adversarial_correction.py} &
Script for adversarial correction using safe, high-only, or high+medium label panels,
including GRL/DAB warm-up options. \\
\filepath{dataset_v3/prepare_real_study_batch_annotation.py} &
Script that parses the uploaded real study-id CSV, aligns it to phase-2 gut samples, and
creates the real-study batch annotation table. \\
\filepath{dataset_v3/real_study_embedding_correction.py} &
Cross-fitted real-study embedding correction baseline using mean-center or mean+scale
correction. \\
\filepath{dataset_v3/outputs_batch_diagnostics_taxonomy/batch_effect_diagnostics_summary.json} &
Main diagnostic summary. \\
\filepath{dataset_v3/outputs_batch_diagnostics_real_study/batch_effect_diagnostics_summary.json} &
Real-study-id diagnostic summary using the uploaded sample-to-study mapping. \\
\filepath{dataset_v3/outputs_batch_diagnostics_taxonomy/probe_metrics.csv} &
Batch, phenotype, and H/D probe metrics. \\
\filepath{dataset_v3/outputs_real_study_embedding_correction/*compact_metrics.csv} &
Compact before/after metrics for cross-fitted real-study embedding correction. \\
\filepath{dataset_v3/outputs_batch_reconstruction_diagnostics_real_study/reconstruction_diagnostics_summary.csv} &
Masked-abundance reconstruction MSE for the base and real-study GRL checkpoints. \\
\filepath{dataset_v3/outputs_batch_diagnostics_taxonomy/*confounding_by_batch.csv} &
Per-batch phenotype-confounding tables. \\
\filepath{dataset_v3/outputs_batch_adversarial_taxonomy/taxonomy_checkpoint_stage2_batch_adversarial_safe.pt} &
Best current conservative adversarial correction checkpoint. \\
\filepath{dataset_v3/outputs_batch_diagnostics_taxonomy/batch_correction_status.md} &
Short status summary for quick reporting. \\
\filepath{dataset_v3/outputs_batch_diagnostics_taxonomy/batch_adversarial_trial_comparison.csv} &
Summary table comparing label-panel ablations and stronger correction runs. \\
\bottomrule
\end{longtable}

\section{Current Conclusion}

The taxonomy-aware BiomeGPT phase-2 embeddings contain batch/study information, so the
batch-correction idea is motivated. However, most candidate batch labels are strongly
phenotype-confounded, so correcting on all batch labels would be scientifically risky.

The earlier proxy-label conservative-safe adversarial trial provided a technical
proof-of-concept: the full pipeline ran, and one light batch-size-8 run reduced safe-batch
predictability with only a small drop in H/D signal. The newer real-study-id
reimplementation is stricter and more informative. With true study ids, the current GRL
fine-tuning preserves or improves masked-abundance reconstruction, but it does not remove
real-study signal. Independent probes show real-study separability increases after the
GRL trials.

The strongest useful result right now is the cross-fitted representation-level correction:
study centering nearly eliminates real-study predictability. But the biological cost is
large when all high-confidence real study labels are used, and much smaller when only
conservative-safe studies are used. Therefore, batch correction is feasible and motivated,
but final model-level correction should be conservative and should optimize against real
study ids, masked-abundance MSE, and H/D preservation together.

\section{Recommended Next Steps}

\begin{enumerate}
  \item Treat the uploaded real study ids as the main batch annotation going forward.
  \item Use conservative-safe real-study labels first; broad high-only study correction is
  too phenotype-confounded.
  \item Add a centroid-alignment/CORAL/MMD-style representation penalty to mimic the
  successful embedding-level study centering inside model training.
  \item Select checkpoints by three criteria: lower real-study probe, preserved or improved
  masked-abundance MSE, and preserved H/D probe.
  \item Fine-tune H/D from the best corrected checkpoint and compare ExVal metrics against
  the uncorrected taxonomy-aware checkpoint.
\end{enumerate}

\end{document}
